% Generated draft: requires mathematical and source review.
\documentclass[11pt]{article}
\usepackage[margin=25mm]{geometry}
\usepackage{fontspec}
\setmainfont{Latin Modern Roman}
\usepackage{amsmath,amssymb,mathtools,unicode-math}
\setmathfont{Latin Modern Math}
\usepackage{enumitem,needspace,xurl,hyperref}
\setlength{\parindent}{0pt}
\setlength{\parskip}{5pt}
\setlength{\emergencystretch}{3em}
\newcommand{\N}{\mathbb{N}}
\newcommand{\Z}{\mathbb{Z}}
\newcommand{\Q}{\mathbb{Q}}
\newcommand{\R}{\mathbb{R}}
\newcommand{\C}{\mathbb{C}}
\newcommand{\F}{\mathbb{F}}
\newcommand{\A}{\mathbb{A}}
\newcommand{\PP}{\mathbb{P}}
\newcommand{\E}{\mathbb{E}}
\newcommand{\eps}{\varepsilon}
\newcommand{\dd}{\,\mathrm{d}}
\begin{document}
\section*{UnsolvedMath ID: 20002000}
\textbf{Source problem code:} catalogue entry
\subsection{Definitions and mathematical statement}
Let $f\in\C[x_1,x_2]$ be nonzero and let $V_f=\{f=0\}\cap(\C^*)^2$. On its smooth locus put
\[
 \operatorname{Log}(z)=(\log|z_1|,\log|z_2|),\qquad
 \gamma_f(z)=[z_1f_{x_1}(z):z_2f_{x_2}(z)]\in\PP^1(\C).
\]
The usual critical locus of $\operatorname{Log}|_{V_f}$ is $\gamma_f^{-1}(\PP^1(\R))$, and its image is the amoeba contour $C_f$. This is distinct from the ramification locus, or critical points, of $\gamma_f$ itself. The source phrase ``locus of the critical points of the Gauss map'' is therefore retained as a definitional ambiguity and is not adopted literally as a subset of $\R^2$.

\textbf{Question A.18, Thorsten Theobald.} When contour singular values are taken as vertices and intervening arcs as edges, what combinatorial restrictions does the planar contour graph have, and which graphs occur for some polynomial $f$? How should the question extend to higher dimensions? The source explains that tracing the contour helps numerically determine the amoeba boundary; the contour may contain interior arcs, so it is not identified with the boundary alone. Degenerate curves and overlaps need conventions before their contours define a finite graph.

\textbf{Supplied stable toric special case.} Assume $V_f$ is smooth and torically nondegenerate with two-dimensional Newton polygon $\Delta$, the compactified logarithmic Gauss map is transverse to $\PP^1(\R)$, and the contour satisfies the explicit stability/end hypotheses of the aid: finite singularities are ordinary cusps or transverse double values, and the ends are recorded individually in their normalization rather than merged solely because their limiting directions or images coincide. Let
\[
 B=\sum_{\text{edges }e\subset\Delta}\operatorname{length}_{\Z}(e)
   =\#(\partial\Delta\cap\Z^2).
\]
The aid asserts there are $B$ toric punctures, giving $2B$ leaves in the decorated normalized-end contour graph, paired by their punctures. Finite vertices have degree two at cusps or inserted markers and degree four at transverse double values. A marker is inserted on an otherwise vertex-free circle so that it is represented as a finite graph.

Writing $n_\times$ for the number of double values, $k$ for the number of graph components and $b_1$ for its first Betti number, the aid's exact combinatorial relation is
\[
 b_1=n_\times-B+k,
 \qquad n_\times=B-k+b_1.
\]
Each graph component has an even number of leaves, and finite odd-valence vertices are excluded in this class. The compactified critical normalization is asserted to be a union of circles covering $\PP^1(\R)$; with positive covering-degree conventions their degrees sum to
\[
 \deg\gamma_f=2\operatorname{Area}(\Delta),
\]
where area gives a fundamental lattice parallelogram area one. Transversality excludes branch points lying over the real projective line. These are necessary restrictions in the specified stable toric class, not sufficiency or a classification for arbitrary polynomials; the leaf count belongs to the decorated normalized-end graph, not an arbitrary image-graph convention.

For a smooth hypersurface in $(\C^*)^N$, $N>2$, one can analogously define the contour as the image of points where $\operatorname{Log}$ has rank less than $N$, equivalently the inverse image of $\PP^{N-1}(\R)$ under the logarithmic Gauss map before taking $\operatorname{Log}$. Its generic image has dimension $N-1$ and requires a stratified-space analogue rather than a planar graph. The source does not supply a higher-dimensional realization conjecture or a complete singularity convention.
\subsection{Short English statement}
Theobald asks which planar graphs arise from amoeba contours and how to extend the question to higher dimensions. The supplied aid gives necessary end, valence and Euler-type restrictions under smooth toric and stability assumptions. The usual contour is formed from logarithm critical values, so the source's Gauss-critical wording and its graph conventions need explicit correction.
\subsection{Sources}
\begin{itemize}
\item Local UnsolvedMath registry, permanent ID 20002000, source code catalogue entry; canonical \texttt{aim-geometry-notes.json}, notes[337].
\item AIM \emph{Amoebas and tropical geometry}, A.18 ``Contour of an amoeba'', contributed by Thorsten Theobald: \url{https://aimath.org/WWN/amoebas/amoebas.pdf}.
\item T.~Theobald, ``Computing amoebas'', \emph{Exp. Math.} 11 (2002), pp.~513--526; M.~Passare and A.~Tsikh, ``Amoebas: their spines and their contours'', preprint (2003), as supplied.
\item Registry research attempt 1, dated 2026-08-14: partial result, rigor ``proved'', novelty confidence low. The machine-generated, non-peer-reviewed aid requires independent verification before citation.
\end{itemize}
\end{document}
